La pregunta de investigación planteada busca determinar el mejor modelo para predecir el volumen de capturas pesqueras en Baja California a partir de variables oceanográficas. Para ello, es pertinente conocer algunos estudios previos donde se hayan aplicado técnicas de modelado y pronóstico similares en contextos pesqueros. A continuación, se presentan trabajos similares que sirvieron de guía y referencia:


Yoon et al. \cite{yoon2020} compararon modelos de máquina de vectores de soporte, bosques aleatorios y redes neuronales profundas para predecir capturas de Caballa, Anchoa y Calamar en zonas de pesca de Corea del Sur, utilizando datos de satélite, oceánicos y climáticos integrados en una cuadrícula geográfica con tamaño unitario de 15kmx15km. Contaban con 4,000 registros de captura entre 2014 y 2016. Los mejores resultados los obtuvieron con DNN, sugiriendo su utilidad para tareas predictivas pesqueras con datos limitados.


Otros autores como Zhao et al.\cite{zhao2021} y Adibi et al.\cite{adibi2020} han explorado técnicas como redes neuronales convolucionales y bosques aleatorios en contextos pesqueros. El trabajo de Zhao et al. se desarrolló en el Mar de China Oriental; utilizaron datos de un Sistema de Monitoreo de Buques para estudiar su comportamiento a través de las trayectorias que siguen al pescar. Su análisis espaciotemporal de las distribuciones del esfuerzo pesquero les ayuda a hacer predicciones a corto plazo; se basan en el descubrimiento de relaciones cronológicas de pesca entre buques para predecir distribuciones, aprovechando las relaciones cronológicas entre buques —el buque que sale primero da la pauta a los buques que salen después—. Usan los comportamientos de las primeras embarcaciones del día como indicadores para las predicciones a corto plazo. A pesar de ser una idea interesante, nosotros no tenemos datos de trayectorias de las embarcaciones, pero podríamos explorar esta posibilidad en el futuro. 

Por otra parte, el trabajo de Adibi et al. fue en el norte del Adriático. Se utilizó una cuadrícula con unidades de 5kmx5km y los datos de desembarque procedían de la lonja de Chioggia. Tomaron en cuenta una distribución ponderada de las capturas, ya que las zonas con más buques pesqueros durante determinados periodos tenían más probabilidades de obtener mayores índices de capturas. Para predecir las capturas por unidad de esfuerzo (CPUE), emplearon el modelo Random Forest, con una previsión ingenua de referencia a modo de comparación. Los resultados mostraron una mejora del 13\% con respecto al modelo de referencia.


Existen también propuestas recientes enfocadas en optimizar arquitecturas de redes recurrentes como LSTM \citep{greff2017,lindemann2021} y de redes de atención \citep{vaswani2023} para procesamiento de series de tiempo, las cuales podrían ser opciones viables dada la naturaleza secuencial de los datos pesqueros; el ensamble de una red recurrente con un modelo de \emph{gradient boosting}, como el que se evalúa en el Capítulo~\ref{cap:resultados}, es una combinación usada en otros dominios de series cortas y ruidosas \citep{shi2022}.


Finalmente, establecer modelos estadísticos autorregresivos como ARIMA y SARIMA como baseline o referencia, ha mostrado ser una estrategia útil en estudios previos como el de Makwinja et al.\cite{makwinja2021} para evaluar el valor predictivo agregado de técnicas más complejas de machine learning y deep learning. El trabajo de Makwinja et al. en el lago Malombe, en Malawi, usa solamente datos de arribos pesqueros, ya que ARIMA es un modelo autorregresivo y para su ajuste no necesita las variables oceanográficas y meteorológicas.


El mecanismo ambiental que resultó decisivo para interpretar estos resultados ---las olas de calor marinas--- se revisa por separado en la Sección~\ref{sec:revision_mhw}.

En conjunto, estos trabajos previos proveen contexto y sustento metodológico para el estudio que se propone. Se evaluarán diversas técnicas de modelado predictivo y se compararán contra métodos estadísticos simples que sirvan de referencia \citep{ho1999,hyndman2021}.
